\documentclass[11pt, a4paper]{article}
\usepackage[utf8]{inputenc}
\usepackage{geometry}
\usepackage{amsmath}
\usepackage{amssymb}
\usepackage{booktabs}
\usepackage{xeCJK} % 如果使用 pdflatex 请注释掉此行并使用其他中文支持方式

\geometry{left=2.5cm, right=2.5cm, top=2.5cm, bottom=2.5cm}

\title{LAB2 关系模式设计初稿与规范化分析}
\author{成员：俞楚凡、赵敬彦}
\date{2026-04-01}

\begin{document}

\maketitle

\section{关系模式初稿 (Relational Schema)}\label{sec:schema}

\textbf{符号说明}：下划线 \underline{字段} 表示主键 (PK)；后缀 $(FK)$ 表示外键。

\subsection{人员与权限模块}
\begin{itemize}
    \item \textbf{人员基表 (People)}: 
    \\ $(\underline{people\_id}, name, gender, phone, email)$
    \item \textbf{学生表 (Student)}: 
    \\ $(\underline{people\_id}(FK), student\_no, grade, major, dep\_id(FK))$
    \item \textbf{教师表 (Teacher)}: 
    \\ $(\underline{people\_id}(FK), staff\_no, title, dept\_id(FK))$
    \item \textbf{用户表 (User)}: 
    \\ $(\underline{user\_id}, people\_id(FK), verification\_status, username, password\_hash, role\_type, dep\_id(FK))$
\end{itemize}

\subsection{空间地理模块}
\begin{itemize}
    \item \textbf{校区表 (Campus)}: 
    \\ $(\underline{campus\_id}, campus\_name, address)$
    \item \textbf{建筑表 (Building)}: 
    \\ $(\underline{building\_id}, building\_name, campus\_id(FK), building\_type, description)$
    \item \textbf{地点表 (Location)}: 
    \\ $(\underline{location\_id}, location\_name, building\_id(FK), facility\_type, description, open\_time)$
\end{itemize}

\subsection{组织、课程与活动模块}
\begin{itemize}
    \item \textbf{院系表 (Department)}: 
    \\ $(\underline{dep\_id}, dep\_name, contact\_info, office\_location\_id(FK), manager\_id(FK), description)$
    \item \textbf{课程表 (Course)}: 
    \\ $(\underline{course\_id}, course\_name, dep\_id(FK), description)$
    \item \textbf{校园活动表 (Event)}: 
    \\ $(\underline{event\_id}, event\_name, event\_type, start\_time, location\_id(FK), host\_dep\_id(FK), description)$
    \item \textbf{教授关系表 (Teaching)}: 
    \\ $(\underline{teacher\_id(FK),\ course\_id(FK)}, semester)$
    \item \textbf{选课关系表 (Enrollment)}: 
    \\ $(\underline{student\_id(FK),\ course\_id(FK)}, semester, grade)$
    \item \textbf{活动参与表 (EventParticipation)}: 
    \\ $(\underline{participant\_id(FK),\ event\_id(FK)}, register\_time)$
\end{itemize}

\subsection{系统日志模块}
\begin{itemize}
    \item \textbf{问答记录表 (QueryRecord)}: 
    \\ $(\underline{record\_id}, user\_id(FK), raw\_question, query\_time, query\_result)$
\end{itemize}

\section{关系设计}

以下四个分块与第~\ref{sec:schema}~节各实体子节一一对应。

\subsection{人员与权限模块}
\begin{itemize}
    \item \textbf{人员 --- 学生}（属于）\quad $1{:}1$\\
    学生是人员的子类型，采用类表继承模式，通过共享主键 $people\_id$ 关联。每位学生恰好对应一条人员基表记录，反之不一定成立。
    \item \textbf{人员 --- 教师}（属于）\quad $1{:}1$\\
    教师是人员的子类型，同样通过 $people\_id$ 继承人员基表。同一 $people\_id$ 不应同时出现在 Student 和 Teacher 中。
    \item \textbf{人员 --- 用户}（注册）\quad $1{:}1$\\
    一个人员至多注册一个系统用户账号，$User.people\_id \rightarrow People.people\_id$。用户表额外承载身份验证与角色权限信息。
\end{itemize}

\subsection{空间地理模块}
\begin{itemize}
    \item \textbf{校区 --- 建筑}（拥有）\quad $1{:}N$\\
    一个校区包含多栋建筑，$Building.campus\_id \rightarrow Campus.campus\_id$。删除校区需级联处理其下所有建筑。
    \item \textbf{建筑 --- 地点}（拥有）\quad $1{:}N$\\
    一栋建筑包含多个地点（教室、食堂、实验室等），$Location.building\_id \rightarrow Building.building\_id$。地理层次为：校区 $\supset$ 建筑 $\supset$ 地点。
\end{itemize}

\subsection{组织、课程与活动模块}
\begin{itemize}
    \item \textbf{院系 --- 学生}（拥有 / 隶属）\quad $1{:}N$\\
    一个院系下有多名学生，$Student.dep\_id \rightarrow Department.dep\_id$。每名学生仅隶属于一个院系。
    \item \textbf{院系 --- 教师}（拥有 / 隶属）\quad $1{:}N$\\
    一个院系下有多名教师，$Teacher.dept\_id \rightarrow Department.dep\_id$。每名教师仅隶属于一个院系。
    \item \textbf{人员 --- 院系}（负责 / 管理）\quad $1{:}1$\\
    每个院系有一名负责人，$Department.manager\_id \rightarrow People.people\_id$。负责人须为本校在册人员。
    \item \textbf{院系 --- 地点}（坐落 / 办公地点）\quad $N{:}1$\\
    院系的办公地点位于某一地点，$Department.office\_location\_id \rightarrow Location.location\_id$。多个院系可共享同一办公地点。
    \item \textbf{院系 --- 课程}（开设）\quad $1{:}N$\\
    一个院系可开设多门课程，$Course.dep\_id \rightarrow Department.dep\_id$。每门课程归属于一个院系。
    \item \textbf{教师 --- 课程}（教授）\quad $M{:}N$\\
    一名教师可教授多门课程，一门课程也可由多名教师共同讲授。通过关联表 $Teaching(\underline{teacher\_id, course\_id})$ 实现，附带 $semester$ 区分学期。
    \item \textbf{学生 --- 课程}（选课）\quad $M{:}N$\\
    一名学生可选修多门课程，一门课程可被多名学生选修。通过关联表 $Enrollment(\underline{student\_id, course\_id})$ 实现，附带 $semester$、$grade$ 记录学期与成绩。
    \item \textbf{学生 --- 校园活动}（参加）\quad $M{:}N$\\
    一名学生可参加多个校园活动，一个校园活动可有多名学生参与。通过关联表 $EventParticipation(\underline{participant\_id, event\_id})$ 实现，附带 $register\_time$ 记录报名时间。
    \item \textbf{院系 --- 校园活动}（举办）\quad $1{:}N$\\
    一个院系可举办多个校园活动，$Event.host\_dep\_id \rightarrow Department.dep\_id$。
    \item \textbf{地点 --- 校园活动}（位于 / 举办地点）\quad $1{:}N$\\
    一个地点可承办多个校园活动，$Event.location\_id \rightarrow Location.location\_id$。活动时间不重叠时同一地点可复用。
\end{itemize}

\subsection{系统日志模块}
\begin{itemize}
    \item \textbf{用户 --- 问答记录}（产生）\quad $1{:}N$\\
    一个用户可产生多条问答记录，$QueryRecord.user\_id \rightarrow User.user\_id$。记录仅追加写入，用于审计与问答质量分析。
\end{itemize}

\vfill

\section{主键与外键设计说明}

\begin{table}[h]
\centering
\begin{tabular}{@{}llll@{}}
\toprule
\textbf{目标表} & \textbf{主键 (PK)} & \textbf{外键 (FK)} & \textbf{设计依据} \\ \midrule
\textbf{Student/Teacher} & $people\_id$ & $people\_id \rightarrow People$ & 类表继承模式，解耦基础信息与角色属性。 \\
\textbf{Building} & $building\_id$ & $campus\_id \rightarrow Campus$ & 建立地理包含关系，支持按校区统计。 \\
\textbf{Location} & $location\_id$ & $building\_id \rightarrow Building$ & 细化建筑内部空间，区分设施类型。 \\
\textbf{Event} & $event\_id$ & $location\_id \rightarrow Location$ & 关联具体地点，支持查询实时动态。 \\ \bottomrule
\end{tabular}
\end{table}

\section{规范化分析 (Normalization Analysis)}

\subsection{数据冗余 (Data Redundancy)}
本设计严格遵循 \textbf{3NF}（第三范式）。通过在 $Course$ 表中仅存储 $teacher\_id$ 而非教师姓名，以及在 $Location$ 表中通过 $building\_id$ 级联，极大地降低了数据空间冗余，保证了信息的单一事实来源。

\subsection{更新异常 (Update Anomalies)}
由于建立了完整的外键约束，若“校区名称”发生变更，只需在 $Campus$ 表中修改一次。由于 $Building$ 和 $Location$ 的层级关系是通过外键维护的，不会出现部分数据未更新导致的空间层级混乱。

\subsection{查询支持情况}
\begin{itemize}
    \item \textbf{复杂查询}：支持多表连接操作。例如查询某校区的所有食堂：
    \[ \text{Result} = \text{Campus} \bowtie \text{Building} \bowtie \sigma_{\text{type='食堂'}}(\text{Location}) \]
    \item \textbf{统计功能}：可直接通过 $Course$ 表的 $dept\_id$ 聚合统计各院系课程数量。
    \item \textbf{NL2SQL 基础}：清晰的实体边界和外键关联为自然语言转 SQL 的语义解析提供了稳定的元数据基础。
\end{itemize}

\end{document}